\subsection*{Independent second reading of the relative inertia-field extension}

Every paragraph of
\texttt{arithmetic\_notes/B8\_relative\_local.tex},
working claims B.META.15 and B.META.16 (assembly identifiers
B.139 and B.140), was read in four consecutive portions on
18 November 2026.  This was a complete reading of the 654-line
version then present, not a check based on its summary.
The outcome is \textbf{PASS}.  The primary comparison inputs
were also independently opened and read at their exact
locators: Temkin, Theorem~1.9.5(ii), \S1.9.7 and
Lemma~1.9.8, printed p.~25; and the complete statement and
proof of Stacks Project, Tag~09EJ, Lemma~15.115.2.
The previously completed full reading of B.136 supplies the
norm-proof comparison used below.  This is a separate
same-model-family check, by GPT-6 Astra Pro in Pro mode;
a separate effort setting is not exposed.  No human,
different-model, or priority claim is made.

\medskip\noindent\textbf{B.139: completion and the chosen algebraic part.}
A finite list of generators of $D/\mathbb Q$ lies in one
finite extension of $\mathbb Q_p$.  The closure of $D$
contains $\mathbb Q_p$ and is therefore itself a finite
local field.  Its restricted valuation on $D$ is discrete,
with finite residue field.  Temkin's exact rank-one
henselization description applies in characteristic zero
and gives $D^h=M_0\cap\Omega$.

The valuation ring of $D^h$ is a DVR and has completion
the valuation ring of $M_0$.  Stacks Tag~09EJ gives
essential surjectivity on finite separable extensions.
For full faithfulness, relative algebraic closedness of
$D^h$ in $M_0$ gives linear disjointness from every finite
Galois extension of $D^h$, hence from every finite
extension.  A finite extension $F/D^h$ is henselian and
relatively algebraically closed in its own completion.
Thus any embedding between completed extensions carries
the algebraic part of one into that of the other.  This
is the required inverse to extending a valued embedding
by continuity; no noncanonical choice of an additional
algebraic closure is being used.

For finite $L/M_0$, a descended extension $F/D^h$ inside
$L$ lies in $\Omega$.  Conversely, an element of
$L\cap\Omega$ is algebraic over $F$ and so belongs to
$F$ by that relative closedness.  Hence $F=L\cap\Omega$,
and its finite field basis is a local basis after
completion.  The same argument applies after every
finite extension of $D$.  The decomposition and inertia
identifications are those for the chosen valuation,
with their respective fixed henselization and strict
henselization fields.  Finally, finite transcendence
degree permits the stated complex embedding; it is
separate from the valuation embedding.

\medskip\noindent\textbf{B.140: the finite-derived Kummer kernel.}
For $r\ge2$, the groups in (META.RL1a) have precisely the
properties required by inflation--restriction.  Namely,
$C=(I^{(r-1)}\cap J)/H$ is abelian and fixes the Tate
coefficient module; each $g\in I^{(r-1)}$ normalizes
$J$, acts trivially on $C$ modulo $H$, and acts trivially
on $\Gamma/C$.  The first assertion uses the commutator
inclusion in $I^{(r)}$, while the second uses normality
of $I^{(r-1)}$ in $I$.  The same $g$ fixes the coefficient
module because $r-1\ge1$.

The two outer terms of the exact sequence are therefore
fixed by $g$.  On its middle term this implies
$(g-1)^2=0$, even though the restriction map need not
be onto the displayed right-hand term.  The element
$a$ has finite orbit in the one finite Galois field
$D_1/D$.  Applying the identity to its Kummer class
gives $m(g-1)\theta=0$ for some positive integer $m$,
hence $(g-1)\theta=0$ over $\mathbb Q_p$.
This does not assert finite order on the whole
cohomology group.

All ratios $g(a)/a$ lie in the single finite local
field $M_1$.  Restriction to its inertia is injective
because $\mathbb Q_p(1)$ has no inertia invariants.
The local rational Kummer kernel is its finite roots
of unity, so one exponent works for every $g$ and
puts a power of $a$ in $F_{D,r-1}$.  The converse is
valid because radicals over that field are abelian,
and it already contains all roots of unity.  This
proves the stable kernel also after finite extensions.
The two-kernel-element identity with $a$ and $1-a$
is elementary and has both implications as stated.

Wild inertia supplies nonzero derived subgroups at
every finite depth through its free pro-$p$ quotients
onto the displayed upper unitriangular groups.
Each nontrivial pro-$p$ group has a nontrivial finite
abelian quotient, so consecutive derived fields are
strict.  A nonzero rational Kummer class remains
nonzero under every finite extension by restriction
and corestriction.  The Gaussian-grid proof uses
the chosen complex embedding only for absolute values
of roots of unity.  The first grid forces $gR=R$;
the translated grid gives a nonzero additive period,
which excludes every nonconstant rational function.

\medskip\noindent\textbf{The nilpotent-union arithmetic input.}
The maximal pronilpotent quotient of inertia has the
stated commuting pro-$p$ and tame pro-prime factors.
The chosen odd prime avoids the local ramification
index as well as $p(p-1)$, so the rational radical
is nontrivial tame and its selected central action
fixes every root of unity.

In the radical extension used for (META.RL3), the
image of the residual kernel $N$ is finite, whereas
the radical Galois group is a torsion-free closed
subgroup of a finite product of $\mathbb Z_p(1)$.
The commutator lies in their intersection and is
therefore trivial.  The relevant subgroup $J_0$
surjects onto that radical Galois group: its fixed
field is exactly the compositum of the finite field,
the relative strict henselization, and the cyclotomic
field used in the proof.  The two restriction
injectivity arguments then reduce to the same
finite-local Kummer kernel.  The positive real
radical and its nontrivial root-of-unity multiple
give the strict absolute-value inequality, a stable
nonzero class, and the two-grid test, including at
$p=2,3$.

\medskip\noindent\textbf{Replacing algebraic numbers by $\Omega$.}
Every finite set of geometric and normalizing data
lies in a finite Galois extension of $D$, whose local
closure remains finite.  Its inertia kernel is
normal in the corresponding local Galois group.
The normalizing Kummer constant is a finite linear
combination of actual values in that local field;
thus its asserted local Galois invariance is justified
without descending a topos morphism.

After choosing the cyclic radical extension, B.139
provides both a basis in $\Omega$ and density of its
algebraic scalar coordinates.  Local equivariance
puts the norm coefficients at those parameters in
the same local field.  The scalar-error argument
uses $[N_1,G_L]\subset N_1\cap G_L$ and triviality of
the Tate action on $N_1$, followed by its surjection
onto $\operatorname{Gal}(L/M_1)$.

The steps from B.136 use algebraicity over the base
field and density, not algebraicity over $\mathbb Q$.
In particular, a retained constant root becomes an
element of $\Omega$ after one actual specialization;
nonzero analytic subresultants can simultaneously
remain nonzero at a dense algebraic point; and
finite interpolation puts the recovered rational
coefficients in $\Omega$.  The two-grid test removes
the common norm factor globally.  All finitely many
coefficients used for graph recovery lie in one
finite extension of the current field, because
$\Omega/E$ is algebraic.  Continuity and henselian
local inverse branches therefore remain available
for pole removal and exceptional-value comparison.

\medskip\noindent\textbf{Infinite algebraic scalar change and descent.}
For $E/F$ algebraic, write $E$ as the filtered union
of its finite subextensions $E_i/F$.  Slicing the
relative topos morphism at the finite etale constant
objects gives compatible morphisms over every $E_i$.
Every finitely presented etale object, morphism,
and finite covering family over $E$ descends to one
$E_i$.  Consequently the compatible slice
inverse-image functors define the limit geometric
morphism over $E$.  This is the precise finite-site
justification of the scalar-change sentence in the
proof.  The relative structural isomorphism is
retained throughout.  Algebraicity of residue
coordinates ensures that closed-point preservation
persists.

The recovered regular function over $E$, viewed
over $\Omega$, has the original $G_F$-equivariant
actual point values.  Equality on all geometric
points gives equality of functions on the reduced
smooth source.  Hence the function is $G_F$-fixed
and descends to $F$.  The same argument after a
finite $F'/F$ uses the finite extension $EF'/E$,
so it verifies the finite-extension quantifier in
the actual-point reconstruction criterion.  In
particular it covers the cases $r=0,1$ by descent
from $r=2$.

\medskip\noindent\textbf{Scope checks.}
The countability argument constructs the claimed
finite transcendence degrees.  Infinitely many
roots of unity exclude embeddings into any
finitely generated extension of a finite local
field: that field's algebraic constant subfield
over the finite local base is finite and has only
finitely many roots of unity.  The ramification
assertion for $r\ge1$ and the nilpotent union is
already forced by the cyclotomic tower.

Every finite quotient of inertia is solvable, so
the union of all finite-derived fields is indeed
$\Omega$.  The semilinear affine-line point map
given at the end is a pinned topos automorphism
over the terminal geometric base but cannot be a
regular polynomial: its values on rational points
would force the polynomial to be the identity.
Thus the proof correctly avoids passing
reconstruction to that infinite union.

\medskip\noindent\textbf{Source and record scope.}
The new source comparison does not rely on a
statement about arbitrary-rank valuations or
inseparable extensions.  The two rank-one,
characteristic-zero inputs were checked directly
at the locators above.  The current proof uses
B.139, B.131, B.135 and B.136 directly; it repeats
the required nilpotent arithmetic argument rather
than treating a stronger field theorem as an
unproved premise.  Exact assembly dependencies
remain the arithmetic agent's responsibility.
This audit was saved separately while the
checkpoint's built proof files and aggregate
ledgers were frozen.
